% theory_qvcg_koszul.tex: Koszul Duality for Quantum Vertex Chiral Groups
% Raeez Lorgat, April 2026

\documentclass[11pt]{article}
\usepackage[localtheorems]{raeez-math-template}

\newtheorem{theorem}{Theorem}[section]
\newtheorem{proposition}[theorem]{Proposition}
\newtheorem{lemma}[theorem]{Lemma}
\newtheorem{corollary}[theorem]{Corollary}
\newtheorem{conjecture}[theorem]{Conjecture}
\theoremstyle{definition}
\newtheorem{definition}[theorem]{Definition}
\newtheorem{example}[theorem]{Example}
\newtheorem{construction}[theorem]{Construction}
\theoremstyle{remark}
\newtheorem{remark}[theorem]{Remark}
\newtheorem{warning}[theorem]{Warning}

\DeclareMathOperator{\Hom}{Hom}
\DeclareMathOperator{\RHom}{RHom}
\DeclareMathOperator{\Ext}{Ext}
\DeclareMathOperator{\HH}{HH}
\DeclareMathOperator{\Rep}{Rep}
\DeclareMathOperator{\Mod}{Mod}
\DeclareMathOperator{\Coh}{Coh}
\DeclareMathOperator{\Fuk}{Fuk}
\DeclareMathOperator{\Perf}{Perf}
\DeclareMathOperator{\Ran}{Ran}
\DeclareMathOperator{\Fact}{Fact}
\DeclareMathOperator{\id}{id}
\DeclareMathOperator{\mult}{mult}
\DeclareMathOperator{\Aut}{Aut}
\DeclareMathOperator{\Tr}{Tr}
\DeclareMathOperator{\rank}{rank}
\DeclareMathOperator{\DT}{DT}
\DeclareMathOperator{\CoHA}{CoHA}
\DeclareMathOperator{\Spec}{Spec}
\DeclareMathOperator{\Op}{Op}
\DeclareMathOperator{\Fun}{Fun}
\DeclareMathOperator{\MC}{MC}
\DeclareMathOperator{\Ob}{Ob}
\DeclareMathOperator{\obs}{obs}

\newcommand{\fkg}{\mathfrak{g}}
\newcommand{\fkh}{\mathfrak{h}}
\newcommand{\fkn}{\mathfrak{n}}
\newcommand{\cA}{\mathcal{A}}
\newcommand{\cC}{\mathcal{C}}
\newcommand{\cH}{\mathcal{H}}
\newcommand{\cM}{\mathcal{M}}
\newcommand{\cO}{\mathcal{O}}
\newcommand{\cR}{\mathcal{R}}
\newcommand{\cW}{\mathcal{W}}
\newcommand{\cZ}{\mathcal{Z}}
\newcommand{\barB}{\overline{B}}
\newcommand{\bC}{\mathbb{C}}
\newcommand{\bH}{\mathbb{H}}
\newcommand{\bZ}{\mathbb{Z}}
\renewcommand{\hbar}{\hslash}

\title{Koszul Duality for Quantum Vertex Chiral Groups}
\author{Raeez Lorgat}
\date{April 2026}

\begin{document}
\maketitle

\begin{abstract}
Koszul duality for quantum vertex chiral groups is a statement about
constructed algebraic layers: the native $E_1$ chiral algebra
$A_X=\Phi_3^{(\Sigma_2,C)}(\cC_X)$, its bar coalgebra $B(A_X)$, the
Koszul dual $A_X^!$, and the center or Drinfeld double carrying the
braided $E_2$ structure. General $G(X)$ is not assumed to exist. The
toric, $K3\times E$, Volume~I, and Hitchin sections below give
conditional root-data prescriptions and explicit proof obligations.
For $\mathbb{C}^3$, the constructed object is the positive half
$Y^+(\widehat{\mathfrak{gl}}_1)$; the full
$\mathcal{W}_{1+\infty}$ layer is recovered only after the
center/double passage.
\end{abstract}

\tableofcontents

% ======================================================================
\section{Quantum vertex chiral groups and Koszul data}
\label{sec:intro}
% ======================================================================

Let $X$ be a Calabi--Yau threefold on a constructed $\Phi_3$ locus. The
specialized output is an $E_1$ chiral algebra
$A_X=\Phi_3^{(\Sigma_2,C)}(\cC_X)$ on the reference curve. A quantum
vertex chiral group $G(X)$ is available only after a positive half,
Hopf or stable-envelope pairing, completion, and center/double
identification have been supplied. When these data exist, its
generalized root datum
$\cR(X) = (\Lambda(X), \Delta^{\mathrm{re}}, \Delta^{\mathrm{im}}_0,
\Delta^{\mathrm{im}}_1, W(X), \rho, \mult)$ encodes the BPS spectrum of
$X$. The bar complex $B(A_X)$ belongs to the native boundary algebra;
the denominator identity belongs to the double or BKM target after the
recognition data are fixed.

\medskip
\noindent\textbf{Four conditional identifications.}
\begin{enumerate}[label=(\arabic*)]
\item For toric CY3 $X_\Sigma$: the positive half is
$Y^+(\widehat{\fkg}_{Q_X})$ on the native $E_1$ side. A Langlands-dual
Yangian statement concerns the completed double/center and requires a
specified Yangian/$\cW$ correspondence (Section~\ref{sec:toric}).

\item For $K3 \times E$: the target BKM denominator is
$64^{-1}\Delta_5(2Z)$ at the $N=1$ Igusa normalization. Its weight is
$\kappa_{\mathrm{BKM}}(\Phi_1)=c_1(0)/2=5$, while the compact
Hodge-supertrace $\kappa_{\mathrm{ch}}(A_{K3\times E})$ is $0$ and the
Heisenberg specialization has $\kappa_{\mathrm{ch}}^{\mathrm{Heis}}=3$
(Section~\ref{sec:k3e}).

\item In the Volume~I language: $A_X$, $B(A_X)$, $A_X^i$, $A_X^!$, and
$Z^{\mathrm{der}}_{\mathrm{ch}}(A_X)$ are separate objects. The
complementarity formula is available only after a cyclic modular
Koszul pair and its conductor have been constructed
(Section~\ref{sec:vol1}).

\item For CY3 arising from the Hitchin system on a curve $C$ with
structure group $\mathbf{G}$: $G(C, \mathbf{G})^! = G(C, {}^L\mathbf{G})$:
geometric Langlands duality is Koszul duality of quantum vertex
chiral groups (Section~\ref{sec:hitchin}).
\end{enumerate}


% ======================================================================
\section{Quantum vertex chiral groups and root data}
\label{sec:recollections}
% ======================================================================

\subsection{The generalized root datum}

On a constructed quantum-vertex-group locus, the generalized root datum
$\cR(X)$ consists of:
\begin{enumerate}[label=(\roman*)]
\item The \emph{lattice} $\Lambda(X) = K_0(\Coh(X))$ (or $K_0(\Fuk(X))$
under HMS) with the Euler form $\chi(\alpha, \beta)$.
\item The \emph{real roots} $\Delta^{\mathrm{re}}$: for $K3 \times E$,
the walls of the fundamental polyhedron $\mathcal{P}_{II}$; for toric
CY3, the vertices of the toric diagram.
\item The \emph{imaginary roots} $\Delta^{\mathrm{im}} =
\Delta^{\mathrm{im}}_0 \sqcup \Delta^{\mathrm{im}}_1$ (even/odd),
with multiplicities $\mult(\alpha) = \DT_\alpha(X)$.
\item The \emph{Weyl group} $W(X)$: generated by reflections in the
real roots.
\item The \emph{Weyl vector} $\rho \in \Lambda(X) \otimes \mathbb{Q}$,
satisfying $(\rho, \delta_i) = -(\delta_i, \delta_i)/2$ for all
simple real roots $\delta_i$.
\end{enumerate}

\subsection{The denominator identity}

When the center/double recognition data identify a BKM target, the
corresponding generalized BKM superalgebra $\fkg_X$ has a
Weyl--Kac--Borcherds denominator identity of product form
\begin{equation}\label{eq:denominator-product}
\Phi_X = e^{-2\pi i \langle \rho, z \rangle}
\prod_{\alpha \in \Delta_+}
\bigl(1 - e^{-2\pi i \langle \alpha, z \rangle}\bigr)^{\mult\,\alpha}.
\end{equation}

For the $N=1$ $K3 \times E$ Borcherds target,
$\Phi_X = 64^{-1}\Delta_5(2Z)$, where $\Delta_5$ is the
theta-normalized Igusa cusp form of weight~5. For toric CY3, the
positive half is governed by the critical CoHA and the DT partition
function supplies product-form evidence; it is not by itself the full
Drinfeld double.

\subsection{The chiral algebra $A_X$ and the bar complex}

The CY-to-chiral assignment on a constructed CY$_3$ locus produces a
native $E_1$ chiral algebra
$A_X = \Phi_3^{(\Sigma_2,C)}(D^b(\Coh(X)))$ with:
\begin{itemize}
\item Bar complex $\barB(A_X)$ quasi-isomorphic to the cyclic bar
complex $\mathrm{CC}_\bullet(\cC)$;
\item Modular characteristic $\kappa_{\mathrm{ch}}(A_X) = \chi^{\mathrm{CY}}(\cC)$;
\item a braided layer only after passage to
$Z^{\mathrm{der}}_{\mathrm{ch}}(A_X)$, to
$\cZ(\Rep^{E_1}(A_X))$, or to a completed Drinfeld double.
\end{itemize}

The product formula~\eqref{eq:denominator-product} is a double or BKM
target identity. The boundary bar coalgebra $B(A_X)$ controls
bar-cobar inversion of $A_X$; the product over positive roots is
identified with it only after a recognition map from the boundary
coalgebra to the denominator target has been constructed.


% ======================================================================
\section{Chiral Koszul duality: review and CY extension}
\label{sec:chiral-koszul-review}
% ======================================================================

\subsection{The chiral Koszul dual}

For a chiral algebra $\cA$ on a smooth curve $C$, the chiral Koszul dual
is formed in two stages:
\[
\cA^i = H^\bullet(\barB(\cA)),\qquad
\cA^! = D_C(\cA^i),
\]
where $D_C$ is Verdier duality on the curve, or the corresponding
linear dual after finite-dimensionality and completion hypotheses are
fixed. The Verdier intertwining gives
\[
D_{\Ran}(\barB(\cA)) \simeq \barB(\cA^!),
\]
identifying the Verdier dual of the bar coalgebra with the bar of the
Koszul dual algebra. Bar-cobar inversion (Theorem~B) gives
$\Omega(\barB(\cA)) \simeq \cA$ on the Koszul locus.

\begin{warning}
Five objects and three functors must be separated:
(0)~$\cA$ (the chiral algebra), $\barB(\cA)$ (the bar coalgebra),
$\cA^i$ (bar cohomology), $\cA^!$ (Verdier dual Koszul algebra), and
$Z^{\mathrm{der}}_{\mathrm{ch}}(\cA)$ (bulk);
(1)~$\Omega(\barB(\cA)) \simeq \cA$ (reconstruction);
(2)~$D_{\Ran}(\barB(\cA)) \simeq \barB(\cA^!)$ (Koszul duality);
(3)~$C^\bullet_{\mathrm{ch}}(\cA, \cA)$ (derived center, bulk observables).
The Koszul dual $\cA^!$ lives on the boundary; the derived center lives
in the bulk.
\end{warning}

\subsection{The complementarity formula}

For a cyclic modular Koszul pair $(\cA, \cA^!)$ with a constructed
conductor $K$, the complementarity formula reads:
\begin{equation}\label{eq:complementarity}
\kappa_{\mathrm{ch}}(\cA) + \kappa_{\mathrm{ch}}(\cA^!) = \varrho \cdot K,
\end{equation}
where $K = c(\cA) + c(\cA^!)$ is the Koszul conductor and
$\varrho$ is the anomaly ratio on that family. The genus-$g$
complementarity is
\[
Q_g(\cA) \oplus Q_g(\cA^!)
\simeq H^*\bigl(\overline{\cM}_g,\, \cZ(\cA)\bigr),
\]
where $\cZ(\cA)$ is the center of $\cA$. Standard examples:
\begin{center}
\renewcommand{\arraystretch}{1.3}
\begin{tabular}{lccccc}
Family & $\cA$ & $\cA^!$ & $K$ & $\kappa_{\mathrm{ch}} + \kappa_{\mathrm{ch}}^!$ \\
\hline
Heisenberg & $H_k$ & $\mathrm{Sym}^{\mathrm{ch}}(V^*)$ & $0$ & $0$ \\
Affine KM & $\widehat{\fkg}_k$ & $\widehat{\fkg}_{-k-2h^\vee}$ & $0$
& $0$ \\
Virasoro & $\mathrm{Vir}_c$ & $\mathrm{Vir}_{26-c}$ & $26$ & $13$ \\
$\cW_N$ & $\cW_N(c)$ & $\cW_N(K_N - c)$ & $K_N$ &
$\varrho_N \cdot K_N$ \\
\end{tabular}
\end{center}
Here $K_N = 4N^3 - 2N - 2$ and $\varrho_N = c \cdot (H_N - 1)/ c =
H_N - 1$ where $H_N = \sum_{i=1}^N 1/i$.

\begin{remark}
For affine Kac--Moody algebras, chiral Koszul duality sends
$k \mapsto -k - 2h^\vee$ (the Feigin--Frenkel involution), and the
complementarity sum $\kappa_{\mathrm{ch}}(\widehat{\fkg}_k) + \kappa_{\mathrm{ch}}(\widehat{\fkg}_{-k-2h^\vee}) = 0$
(this vanishing is specific to KM and free fields; for Virasoro,
$\kappa_{\mathrm{ch}} + \kappa_{\mathrm{ch}}^! = 13$, and for $\cW_N$,
$\kappa_{\mathrm{ch}} + \kappa_{\mathrm{ch}}^! = \varrho_N \cdot K_N \neq 0$ in general).
At critical level
$k = -h^\vee$, the bar complex is uncurved ($\kappa_{\mathrm{ch}} = 0$) and the
Koszul dual is $\widehat{\fkg}_{-h^\vee}$ itself. The Feigin--Frenkel
center $\mathfrak{z}(\widehat{\fkg}_{-h^\vee}) \simeq
\Fun(\Op_{\fkg^\vee}(D^\times))$ connects the critical-level Koszul
dual to the Langlands dual Lie algebra $\fkg^\vee$.
This connection is external motivation. In the CY$_3$ setting it becomes
a theorem only after the positive half, center/double, and
Verdier-duality data have been specified.
\end{remark}


\subsection{Extension to quantum vertex chiral groups}

\begin{definition}[Koszul dual root datum]
\label{def:koszul-dual-root-datum}
Given a generalized root datum $\cR(X) = (\Lambda, \Delta^{\mathrm{re}},
\Delta^{\mathrm{im}}_0, \Delta^{\mathrm{im}}_1, W, \rho, \mult)$,
the \emph{Koszul dual root datum} $\cR(X)^!$ is defined by:
\begin{enumerate}[label=(\roman*)]
\item The lattice $\Lambda^! = \Lambda^\vee$ (the dual lattice with
the negated Euler form $-\chi$).
\item The real roots $(\Delta^{\mathrm{re}})^!$: unchanged as a set,
but with the negated Gram matrix (up to conjugation by $W$).
\item The imaginary root multiplicities: $\mult^!(\alpha) =
\mult(\omega_B(\alpha))$, where $\omega_B$ is the Borcherds involution
(Definition~\ref{def:borcherds-involution}).
\item The Weyl group: $W^! = W$ (unchanged).
\item The Weyl vector: $\rho^! = -\rho + \delta_K$, where $\delta_K$
is the Koszul conductor shift (determined by~\eqref{eq:complementarity}).
\end{enumerate}
\end{definition}

\begin{warning}
The Koszul dual root datum is \emph{not} obtained by simply negating
the Cartan matrix $A \mapsto -A$. The real roots undergo a sign change
in the bilinear form, but the imaginary roots are permuted by the
Borcherds involution $\omega_B$, which acts nontrivially on the
positive cone. The distinction is essential for BKM superalgebras:
negating the Cartan matrix produces an algebra with the wrong root
multiplicities.
\end{warning}


% ======================================================================
\section{Toric CY3: the dual Yangian and $\cW$-algebra correspondence}
\label{sec:toric}
% ======================================================================

\subsection{Setup}

Let $X_\Sigma$ be a toric CY3 determined by a fan $\Sigma$ with
associated quiver $(Q_X, W_X)$. The native Hall object is the positive
half of the affine super Yangian:
\begin{itemize}
\item Positive half: the critical CoHA
$\cH(Q_X, W_X) \simeq Y^+(\widehat{\fkg}_{Q_X})$;
\item $R$-matrix: a structure on the Drinfeld double
$D(Y^+(\widehat{\fkg}_{Q_X}))$ after a Hopf pairing and completion are
fixed;
\item Root multiplicities: DT invariants $\DT_\mathbf{d}(X_\Sigma)$.
\end{itemize}

\subsection{The Yangian--$\cW$-algebra duality}

The required comparison is the Yangian/$\cW$-algebra correspondence
(Maulik--Okounkov, Schiffmann--Vasserot, Prochazka--Rapcak):
\begin{equation}\label{eq:yangian-w-duality}
Y(\widehat{\fkg}) \longleftrightarrow \cW(\fkg^L),
\end{equation}
where $\fkg^L$ is the Langlands dual of the super Lie algebra $\fkg$
attached to the quiver. The precise statement: the category of
representations of the affine Yangian $Y(\widehat{\fkg})$ at a
particular level is equivalent, as a braided monoidal category, to
the category of modules over the $\cW$-algebra $\cW(\fkg^L)$ at the
dual level.

\begin{proposition}[Toric Koszul-dual proof obligation]
\label{prop:toric-koszul-dual}
For a toric CY3 $X_\Sigma$ without compact $4$-cycles, the assertion
\begin{equation}\label{eq:toric-koszul}
G(X_\Sigma)^! \simeq Y(\widehat{\fkg^L}_{Q_X}),
\end{equation}
for the completed quantum vertex group follows only from the following
data: a constructed positive half $Y^+(\widehat{\fkg}_{Q_X})$, a
non-degenerate Hopf pairing, a completed Drinfeld double, and a
Yangian/$\cW$ correspondence compatible with Koszul duality. The target
diagram is
\[
Y(\widehat{\fkg}_{Q_X}) \xleftrightarrow{\;\text{Y/W}\;}
\cW(\fkg^L_{Q_X}) \xleftarrow{\;\text{Koszul}\;}
Y(\widehat{\fkg^L}_{Q_X}).
\]
\end{proposition}

\begin{proof}[Argument]
The native chiral algebra $A_{X_\Sigma}$ is $E_1$. Its positive Hall
comparison is $Y^+(\widehat{\fkg}_{Q_X})$. After a pairing and
completion are supplied, the double has a braided representation
category controlled by the Yangian $R$-matrix. The Yangian/$\cW$
correspondence identifies this braided category with
$\Rep(\cW(\fkg^L_{Q_X}))$ at an appropriate level.

The $\cW$-algebra Koszul duality (Volume~I, with conductor
$K_N = 4N^3 - 2N - 2$ for $\cW_N$) sends
$\cW_k(\fkg^L) \mapsto \cW_{k^!}(\fkg^L)$, where $k^!$ is
determined by the complementarity formula. For the $\cW_{1+\infty}$
algebra (the $\bC^3$ double/center case), this is the self-dual point
where the Yangian/$\cW$ correspondence identifies the completed double
of $Y^+(\widehat{\mathfrak{gl}}_1)$ with $\cW_{1+\infty}$ at the
self-dual level. Applying the Yangian/$\cW$ correspondence in reverse on the dual
side gives $\cW_{k^!}(\fkg^L) \leftrightarrow
Y(\widehat{\fkg^L}_{Q_X})$, yielding the
conditional identification~\eqref{eq:toric-koszul}.
\end{proof}

\begin{example}[$\bC^3$: self-dual Yangian]
\label{ex:c3-koszul}
For $X = \bC^3$, the constructed Hall object is
$Y^+(\widehat{\mathfrak{gl}}_1)$. The full $\cW_{1+\infty}$ algebra
appears after the Drinfeld double/center passage. The super Lie algebra is
$\fkg = \mathfrak{gl}_1$, which is self-dual:
$\mathfrak{gl}_1^L = \mathfrak{gl}_1$. Hence
\[
D(Y^+(\widehat{\mathfrak{gl}}_1))^!
\simeq D(Y^+(\widehat{\mathfrak{gl}}_1))
\]
after the completed Hopf pairing is fixed. The $\bC^3$ double is
Koszul self-dual.
This is consistent with the self-duality of $\cW_{1+\infty}$ at the
self-dual level (Volume~I: the Yangian $R$-matrix is the DK-0 shadow,
and MC4 is solved by weight stabilization).
\end{example}

\begin{example}[Resolved conifold]
\label{ex:conifold-koszul}
For the resolved conifold $X = \cO(-1) \oplus \cO(-1) \to \bP^1$:
the quiver is Klebanov--Witten (two vertices, arrows in both
directions), and $\fkg_{Q_X} = \mathfrak{gl}_{1|1}$, the general
linear superalgebra. Since $\mathfrak{gl}_{1|1}$ is self-dual
under the Langlands involution (it has a single even and single odd
simple root, and the Dynkin diagram is symmetric), we obtain
\[
G(X_{\mathrm{conifold}})^! \simeq
Y(\widehat{\mathfrak{gl}}_{1|1}) \simeq G(X_{\mathrm{conifold}}).
\]
This is a double-level self-duality claim. The native conifold
$\Phi_3$ output remains an $E_1$ positive-half object.
The self-duality reflects the geometric involution
$\cO(-1) \oplus \cO(-1) \to \cO(-1) \oplus \cO(-1)$ exchanging the
two line bundles.
\end{example}

\subsection{The negative mode algebra interpretation}

In the Yangian literature, the ``dual'' of the positive modes
$Y^+(\widehat{\fkg})$ (generators $e_i(u)$, $\psi_i(u)$) consists of
the negative modes $Y^-(\widehat{\fkg})$ (generators $f_i(u)$,
$\psi_i(u)^{-1}$). The Koszul dual root datum has $\rho^! = -\rho$
(reflecting the sign flip on the positive cone), and the bar complex
$\barB(Y^+) \simeq Y^-$ encodes the passage from creation to
annihilation operators.

This is the $E_1$ boundary sector. The braided Koszul statement
operates after the center/double passage and incorporates the
Langlands involution on the root datum, exchanging $\fkg$ with
$\fkg^L$. The negative mode
algebra is $Y^-(\widehat{\fkg})$; the Koszul dual quantum vertex
chiral group is $Y(\widehat{\fkg^L})$, the full algebra of the
\emph{Langlands-dual} super Lie algebra.

\subsection{Root multiplicities and the dual DT partition function}

The Koszul dual root multiplicities are determined by
\begin{equation}\label{eq:dual-dt}
\mult^!(\alpha) = \DT_{\omega_B(\alpha)}(X_\Sigma),
\end{equation}
where $\omega_B$ is the Borcherds involution on the positive cone of
$\Lambda(X_\Sigma)$. For toric CY3 with the standard toric involution
$\sigma \colon X_\Sigma \to X_\Sigma$ (induced by the involution of the
fan), $\omega_B$ agrees with $\sigma^*$ on $K_0$.

The dual DT partition function is
\[
Z^{\DT}(X_\Sigma)^! =
\prod_{\alpha \in \Delta_+}
\bigl(1 - e^{-2\pi i \langle \alpha, z \rangle}\bigr)^{%
\DT_{\omega_B(\alpha)}(X_\Sigma)}.
\]


% ======================================================================
\section{$K3 \times E$: mirror symmetry and the Borcherds involution}
\label{sec:k3e}
% ======================================================================

\subsection{The BKM superalgebra $\fkg_{\Delta_5}$ and its Koszul dual}

At $N=1$, the $K3\times E$ Borcherds target is the generalized BKM
superalgebra $\fkg_{\Delta_5}$. This is target BKM data; a compact
Hall realization additionally requires the finite Hall--Borcherds
recognition datum. The target has:
\begin{itemize}
\item Real simple roots $\{\delta_1, \delta_2, \delta_3\}$ with Gram
matrix
$(\delta_i, \delta_j) = \begin{psmallmatrix} 2 & -2 & -2 \\ -2 & 2 & -2 \\
-2 & -2 & 2 \end{psmallmatrix}$;
\item Imaginary root multiplicities $\mult(\alpha) = f(nm, l)$,
the Fourier coefficients of $\phi_{0,1}$ (the K3 elliptic genus);
\item Denominator identity:
$\operatorname{den}(\fkg_{\Delta_5})=64^{-1}\Delta_5(2Z)$.
\end{itemize}

\begin{definition}[Borcherds involution]
\label{def:borcherds-involution}
The \emph{Borcherds involution} $\omega_B$ on $\fkg_{\Delta_5}$ is the
involution determined by:
\begin{enumerate}[label=(\roman*)]
\item On real roots: $\omega_B(\delta_i) = -\delta_i$ (the Chevalley
involution $e_i \mapsto -f_i$, $f_i \mapsto -e_i$,
$h_i \mapsto -h_i$).
\item On imaginary roots: $\omega_B$ permutes the imaginary simple
roots within the positive cone $\cR_{\geq 0} \mathcal{P}_{II}$
according to the involution
\begin{equation}\label{eq:borcherds-imaginary}
\omega_B \colon (n, l, m) \longmapsto (m, -l, n)
\end{equation}
on the lattice coordinates, reflecting the Fourier expansion through
$f(nm, l) = f(nm, -l)$ (the symmetry $l \mapsto -l$ of
$\phi_{0,1}$) composed with the transposition $n \leftrightarrow m$.
\end{enumerate}
\end{definition}

\begin{warning}
The Borcherds involution is \emph{not} the naive negation
$\alpha \mapsto -\alpha$ on all roots. On real roots it agrees with
negation (up to Weyl conjugacy), but on imaginary roots it acts as
the specific permutation~\eqref{eq:borcherds-imaginary}. The
distinction is essential: the imaginary root multiplicities $f(nm, l)$
are not invariant under arbitrary permutations of the lattice
coordinates, only under $\omega_B$.
\end{warning}

\subsection{The mirror CY3 and modular transformation}

\begin{proposition}[Borcherds-target dual denominator datum]
\label{prop:k3e-koszul-denominator}
The dual denominator identity of the $K3\times E$ BKM target is
controlled by the automorphic transformation law of the corresponding
Siegel modular form. At $N=1$,
\[
 \operatorname{den}(\fkg_{\Delta_5})=64^{-1}\Delta_5(2Z),
 \qquad
 \kappa_{\mathrm{BKM}}(\Phi_1)=\mathrm{wt}(\Delta_5)=c_1(0)/2=5 .
\]
For a CHL level $N$, replacing $\Delta_5$ by the correct row form
$\Phi_N$ gives $\kappa_{\mathrm{BKM}}(\Phi_N)=c_N(0)/2$. A Fricke
formula must use the actual multiplier, the factor $\det(CZ+D)$, and
the row's weight. It is not a formula with
$\kappa_{\mathrm{ch}}(A_X)$ in the exponent.
\end{proposition}

\begin{proof}[Sketch]
The denominator identity is target arithmetic. Igusa normalization gives
$D_5=64^{-1}\Delta_5$, hence the denominator
$64^{-1}\Delta_5(2Z)$. The modular exponent is the BKM weight,
not the compact chiral Hodge supertrace: for compact $K3\times E$,
$\kappa_{\mathrm{ch}}(A_{K3\times E})=0$ by Serre cancellation, while
the Heisenberg specialization gives $\kappa_{\mathrm{ch}}^{\mathrm{Heis}}=3$.
The equality $5=c_1(0)/2$ is the Borcherds-weight identity.
\end{proof}

\subsection{Geometric interpretation: mirror symmetry}

The Fricke involution on $\bH_2$ is expected to correspond, on the
geometric side, to the passage from $X = (S \times E)/(\bZ/N\bZ)$ to its
\emph{mirror} CY3:
\[
X^\vee = (S^\vee \times E^\vee) / (\bZ/N\bZ),
\]
where $S^\vee$ is the mirror K3 (determined by the mirror map on
the K3 moduli space) and $E^\vee$ is the dual elliptic curve
$E^\vee = \bC / (\bZ + N\tau \bZ)$ (the $N$-isogeny dual).

\begin{conjecture}[Mirror Koszul duality for $K3 \times E$]
\label{conj:k3e-mirror-koszul}
The Koszul dual quantum vertex chiral group of $K3 \times E$ is
the quantum vertex chiral group of the mirror:
\[
G(X)^! \simeq G(X^\vee).
\]
Under this identification:
\begin{enumerate}[label=(\roman*)]
\item The Borcherds involution $\omega_B$ on imaginary roots
corresponds to the mirror map on the BPS spectrum:
$\DT_\alpha(X^\vee) = \DT_{\omega_B(\alpha)}(X)$.
\item The Koszul conductor must be computed from the cyclic modular
Koszul pair. It is not determined by replacing the Borcherds weight
with $\kappa_{\mathrm{ch}}(A_X)$.
\item The eight diagonal-divisor Siegel forms give eight BKM
denominator targets. Compact $K3\times E$ CHL hosts occur only on the
appropriate CHL subrows; the remaining rows are automorphic targets
until a host geometry and recognition map are supplied.
\end{enumerate}
\end{conjecture}

\begin{remark}[Self-mirror and self-duality]
At the self-mirror point $X \simeq X^\vee$ (the attractor point in
K3 moduli), the BKM superalgebra $\fkg_{\Delta_5}$ becomes Koszul
self-dual. This is the CY3 analogue of the Virasoro self-duality at
$c = 13$: the self-dual datum is the mirror-fixed Borcherds target
together with its multiplier and weight, not a replacement of
$\kappa_{\mathrm{BKM}}$ by $\kappa_{\mathrm{ch}}$.
\end{remark}

\subsection{The Borcherds involution is not Cartan negation}

The following example illustrates why the Koszul dual of a BKM
superalgebra cannot be obtained by negating the Cartan matrix.

\begin{example}
Consider the real root Gram matrix
$A = \begin{psmallmatrix} 2 & -2 & -2 \\ -2 & 2 & -2 \\
-2 & -2 & 2 \end{psmallmatrix}$.
The naive ``Cartan-negated'' BKM algebra $\fkg_{-A}$ has Gram
matrix $-A = \begin{psmallmatrix} -2 & 2 & 2 \\ 2 & -2 & 2 \\
2 & 2 & -2 \end{psmallmatrix}$, which has signature $(1, 2)$ rather
than $(2, 1)$. The resulting algebra has \emph{no} positive-definite
Cartan subalgebra and is not a well-defined BKM algebra in the sense
of Borcherds.

Even if one formally defines $\fkg_{-A}$, its root
multiplicities would be determined by a different automorphic form
(one whose Weyl denominator involves $-A$), not by the Fourier
coefficients of $\phi_{0,1}$ permuted by $\omega_B$. The Koszul dual
$\fkg_{\Delta_5}^!$ has the \emph{same} real root Gram matrix $A$
(up to Weyl conjugacy) but \emph{permuted} imaginary root
multiplicities. The permutation is $\omega_B$, not the identity or
negation.
\end{example}


% ======================================================================
\section{Volume~I framework: $\kappa_{\mathrm{ch}}(A^!_X)$ and complementarity}
\label{sec:vol1}
% ======================================================================

\subsection{The CY complementarity formula}

\begin{proposition}[CY complementarity datum]
\label{thm:cy-complementarity}
For a CY$_3$ constructed locus with native $E_1$ chiral algebra $A_X$,
a Koszul dual $A^!_X$, and a cyclic modular Koszul conductor $K_X$, the
expected complementarity relation is
\begin{equation}\label{eq:cy-complementarity}
\kappa_{\mathrm{ch}}(A_X) + \kappa_{\mathrm{ch}}(A^!_X) = \varrho_X \cdot K_X,
\end{equation}
where:
\begin{itemize}
\item $K_X$ is the Koszul conductor of the pair $(A_X,A_X^!)$;
\item $\varrho_X = \kappa_{\mathrm{ch}}(A_X) / c_{\mathrm{eff}}(X)$ is the anomaly
ratio, where $c_{\mathrm{eff}}$ is the effective central charge of
$A_X$.
\end{itemize}
The proof obligation is to construct $A_X^!$, $K_X$, and the genus-one
trace identifying $\kappa_{\mathrm{ch}}(A_X)$ with the appropriate
negative-cyclic Hodge summand. This is the CY extension of the
$\cW$-algebra complementarity formula
$\kappa_{\mathrm{ch}} + \kappa_{\mathrm{ch}}^! = \varrho_N \cdot K_N$.
\end{proposition}

\begin{proof}[Criterion]
The genus-$g$ complementarity (Theorem~C of Volume~I) gives
\[
Q_g(A_X) \oplus Q_g(A^!_X) \simeq
H^*\bigl(\overline{\cM}_g,\, \cZ(A_X)\bigr).
\]
Taking the leading Hodge class $\lambda_g$ contribution at genus~$g$:
$\obs_g(A_X) + \obs_g(A^!_X) = \kappa_{\mathrm{ch}}(A_X) \cdot \lambda_g +
\kappa_{\mathrm{ch}}(A^!_X) \cdot \lambda_g = (\kappa_{\mathrm{ch}} + \kappa_{\mathrm{ch}}^!) \cdot \lambda_g$.
The right-hand side is determined by the center
$Z^{\mathrm{der}}_{\mathrm{ch}}(A_X)$ after the center exists. The
displayed formula follows precisely when its genus-one Hodge summand is
computed and equals $\varrho_X K_X$.
\end{proof}

\subsection{The shadow obstruction tower of $A^!_X$}

The shadow obstruction tower of the Koszul dual $A^!_X$ is determined by the
shadow obstruction tower of $A_X$ via the Koszul intertwining:

\begin{proposition}[Shadow obstruction tower duality]
\label{prop:shadow-duality}
After the cyclic modular Koszul pair and conductor are constructed, the
MC elements $\Theta_{A_X}$ and $\Theta_{A^!_X}$ should be related by
\begin{equation}\label{eq:shadow-duality}
\Theta_{A^!_X} = -\Theta_{A_X} + \Theta^{\mathrm{conductor}}_X,
\end{equation}
where $\Theta^{\mathrm{conductor}}_X$ is the conductor contribution
determined by $K_X$. At degree $2$ (the scalar level):
$\kappa_{\mathrm{ch}}(A^!_X) = -\kappa_{\mathrm{ch}}(A_X) + \varrho_X \cdot K_X$,
recovering~\eqref{eq:cy-complementarity}.
\end{proposition}

At the level of denominator identities, equation~\eqref{eq:shadow-duality}
gives the relation between $\Phi_X$ and $\Phi_{X^!}$: the product
over positive roots of $G(X)^!$ is obtained from the product over
positive roots of $G(X)$ by the Borcherds involution on the
imaginary roots plus the conductor shift on the Weyl vector.

\subsection{Classification by shadow depth}

The shadow depth classification (Volume~I: classes G/L/C/M) extends to
constructed quantum vertex chiral groups. Preservation under Koszul
duality is a criterion, not a formal consequence of the bar-cobar
counit:

\begin{proposition}[Shadow depth preservation]
\label{prop:depth-preservation}
If $G(X)$ and $G(X)^!$ are constructed and the Koszul involution
preserves the quartic shadow invariant $S_4$, then the two have the
same shadow depth $r_{\max}$.
\end{proposition}

\begin{proof}
Shadow depth is determined by the critical discriminant
$\Delta = 8\kappa_{\mathrm{ch}} S_4$. Under the stated hypothesis
$S_4^! = S_4$. The discriminant complementarity formula
$\Delta(A) + \Delta(A^!) = \frac{6960}{(5c+22)(152-5c)}$ ensures
that $\Delta = 0$ if and only if $\Delta^! = 0$ (at the self-dual
point $c_* = 13$ for Virasoro; the analogous statement holds for
CY3 quantum vertex chiral groups only after their center/double model
and conductor are supplied).
\end{proof}


% ======================================================================
\section{The Hitchin system: Langlands duality as Koszul duality}
\label{sec:hitchin}
% ======================================================================

\subsection{CY3 from the Hitchin system}

Let $C$ be a smooth projective curve of genus $g_C \geq 2$ and
$\mathbf{G}$ a connected reductive group. The \emph{Hitchin system}
$\cM_H(C, \mathbf{G}) = T^*\mathrm{Bun}_\mathbf{G}(C)$ is a
holomorphic symplectic manifold (an algebraic completely integrable
system). The total space of the Hitchin fibration
\[
h \colon \cM_H(C, \mathbf{G}) \longrightarrow
\cA_H = \bigoplus_{i=1}^{\rank(\mathbf{G})}
H^0(C, \omega_C^{\otimes (d_i + 1)})
\]
is a (non-compact) holomorphic symplectic variety of dimension
$2 \dim \mathrm{Bun}_\mathbf{G}(C) = 2(g_C - 1) \dim \mathbf{G}$.

When $\dim_\bC \cM_H = 6$, i.e., when
$(g_C - 1) \dim \mathbf{G} = 3$, the Hitchin moduli space is a
CY3 (more precisely, a non-compact CY3 with trivial canonical bundle).
This occurs for:
\begin{itemize}
\item $\mathbf{G} = \mathrm{SL}_2$, $g_C = 2$:
$\dim \cM_H = 2 \cdot 1 \cdot 3 = 6$.
\item $\mathbf{G} = \mathrm{SL}_3$, $g_C = 2$:
$\dim \cM_H = 2 \cdot 1 \cdot 8 = 16$ (not CY3).
\item More generally, for any $\mathbf{G}$ and $g_C$ with the right
dimension.
\end{itemize}

Even when $\dim \cM_H \neq 6$, the Hitchin system
carries a \emph{CY category}: the derived category
$D^b(\Coh(\cM_H(C, \mathbf{G})))$ is a CY category of dimension
$\dim \cM_H$, and the Fukaya category $\Fuk(\cM_H(C, \mathbf{G}))$
is a CY category of the same dimension. The CY-to-chiral functor
$\Phi_d$ can therefore be applied only in its dimension-stratified
form. For $d\geq 3$ the specialized output is $E_1$-chiral; a quantum
vertex group $G(C,\mathbf{G})$ requires the same positive-half,
pairing, completion, and center/double data as in the CY$_3$ examples.

\subsection{Langlands duality of Hitchin systems}

The geometric Langlands duality of Hitchin systems (Hausel--Thaddeus,
Donagi--Pantev) gives a canonical identification
\begin{equation}\label{eq:hitchin-langlands}
\cM_H(C, \mathbf{G})^\vee \simeq \cM_H(C, {}^L\mathbf{G}),
\end{equation}
where ${}^L\mathbf{G}$ is the Langlands dual group. More precisely:
\begin{enumerate}[label=(\roman*)]
\item The Hitchin bases are canonically identified:
$\cA_H(\mathbf{G}) \simeq \cA_H({}^L\mathbf{G})$ (both are determined
by the invariant polynomials, which depend only on the root datum);
\item Generic fibers are dual abelian varieties: for a regular
semisimple Hitchin fiber $h^{-1}(a) \simeq \mathrm{Prym}_{C,\mathbf{G}}(a)$,
the dual is $h^{-1}(a)^\vee \simeq \mathrm{Prym}_{C,{}^L\mathbf{G}}(a)$;
\item At the level of derived categories (Homological Mirror Symmetry
for Hitchin systems):
$D^b(\Coh(\cM_H(\mathbf{G}))) \simeq \Fuk(\cM_H({}^L\mathbf{G}))$.
\end{enumerate}

\subsection{The key conjecture}

\begin{conjecture}[Langlands duality is Koszul duality]
\label{conj:langlands-koszul}
For a curve $C$ and reductive group $\mathbf{G}$, the Koszul dual
of the quantum vertex chiral group $G(C, \mathbf{G})$ is the
quantum vertex chiral group of the Langlands-dual group:
\begin{equation}\label{eq:langlands-koszul}
\boxed{G(C, \mathbf{G})^! \simeq G(C, {}^L\mathbf{G}).}
\end{equation}
\end{conjecture}

\begin{remark}[Scope]
This is a conjecture about the \emph{quantum vertex chiral groups}
$G(C, \mathbf{G})$, not about the chiral algebras of the Hitchin system.
The passage from the CY category $D^b(\Coh(\cM_H))$ to the quantum
vertex chiral group involves the full CY-to-chiral functor $\Phi$
(Theorem~CY-A) and the root datum extraction. The conjecture asserts
that Koszul duality, when expressed at the level of root data,
coincides with the Langlands involution on the root datum of
$\mathbf{G}$.
\end{remark}

\subsection{Evidence}

The evidence for Conjecture~\ref{conj:langlands-koszul} comes from
four independent directions.

\subsubsection{(a) Root datum level}
The Langlands dual group ${}^L\mathbf{G}$ has root datum
$(\Lambda^\vee, \Delta^\vee, \Delta, W)$: the coroots become
roots, the cocharacter lattice becomes the character lattice, and the
Weyl group is unchanged. The Koszul dual root datum
(Definition~\ref{def:koszul-dual-root-datum}) has
$\Lambda^! = \Lambda^\vee$ and $W^! = W$. The real root systems agree.

For simply-laced $\mathbf{G}$ (types $A$, $D$, $E$): $\fkg^L = \fkg$,
so the conjecture predicts Koszul self-duality
$G(C, \mathbf{G})^! \simeq G(C, \mathbf{G})$, with the self-dual
point at a specific level/central charge.

For non-simply-laced $\mathbf{G}$ (types $B$, $C$, $F_4$, $G_2$):
the Langlands involution exchanges long and short roots, and the
conjecture predicts a genuine duality between different quantum vertex
chiral groups.

\subsubsection{(b) Chiral algebra level: Feigin--Frenkel at critical level}

At critical level $k = -h^\vee$, the affine algebra
$\widehat{\fkg}_{-h^\vee}$ has:
\begin{itemize}
\item $\kappa_{\mathrm{ch}}(\widehat{\fkg}_{-h^\vee}) = 0$ (uncurved bar complex);
\item Center $\mathfrak{z}(\widehat{\fkg}_{-h^\vee}) \simeq
\Fun(\Op_{\fkg^\vee}(D^\times))$ (Feigin--Frenkel);
\item Critical-level complementarity:
$Q_g(\widehat{\fkg}_{-h^\vee}) \oplus Q_g(\widehat{\fkg}_{-h^\vee})
\simeq H^*(\overline{\cM}_g, \Fun(\Op_{\fkg^\vee}))$.
\end{itemize}

The Feigin--Frenkel center provides the link: the critical-level
chiral Koszul dual has its center controlled by
$\Op_{\fkg^\vee}$, the opers for the Langlands dual Lie algebra.
In the quantum vertex chiral group setting, this becomes
$G(C, \mathbf{G})_{k=-h^\vee}^! \simeq G(C, {}^L\mathbf{G})_{k^L = -h^{\vee,L}}$
at critical levels.

\subsubsection{(c) Shifted Yangian level: symplectic duality}

The shifted Yangian conjecture (Volume~I, Conjecture~26.3.2):
\begin{equation}\label{eq:shifted-yangian-langlands}
Y_\mu(\fkg)^! \simeq Y_{-\mu}(\fkg^\vee),
\end{equation}
is an $E_1$-chiral incarnation of symplectic duality
(Braden--Licata--Proudfoot--Webster). For unshifted $\mu = 0$ and
simply-laced $\fkg$, this reduces to Yangian self-duality.

The quantum vertex chiral group $G(C, \mathbf{G})$ is built from the
CoHA of the Hitchin quiver, which is a shifted Yangian. The conjecture
$G(C, \mathbf{G})^! = G(C, {}^L\mathbf{G})$ extends the shifted
Yangian duality~\eqref{eq:shifted-yangian-langlands} to the full
center/double setting.

\subsubsection{(d) Homological mirror symmetry}

HMS for Hitchin systems~\eqref{eq:hitchin-langlands} gives
$D^b(\Coh(\cM_H(\mathbf{G}))) \simeq \Fuk(\cM_H({}^L\mathbf{G}))$.
Applying $\Phi_d$ gives an identification of chiral outputs only when
the HMS equivalence preserves the CY structure, the specialization
datum, and the obstruction-killing package:
\[
A_{C,\mathbf{G}} = \Phi_d(D^b(\Coh(\cM_H(\mathbf{G}))))
\simeq \Phi_d(\Fuk(\cM_H({}^L\mathbf{G})))
= A_{C,{}^L\mathbf{G}}.
\]
But this is the \emph{HMS identification}, not Koszul duality. The
relationship between HMS and Koszul duality is:
\[
\text{HMS: } A_{C,\mathbf{G}} \simeq A_{C,{}^L\mathbf{G}}
\quad\text{vs.}\quad
\text{Koszul: } A_{C,\mathbf{G}}^! \simeq A_{C,{}^L\mathbf{G}}.
\]
These are consistent if and only if $A_{C,\mathbf{G}}$ is Koszul
self-dual (for simply-laced $\mathbf{G}$). For non-simply-laced
$\mathbf{G}$, HMS gives $A_{C,\mathbf{G}} \simeq A_{C,{}^L\mathbf{G}}$
as underlying chiral algebras, but the Koszul involution acts
nontrivially on the center-layer braiding, exchanging the quantum group
parameter $q$ and $q^{-1}$.

\subsection{The Koszul conductor of the Hitchin CY3}

\begin{proposition}[Hitchin conductor proof obligation]
\label{prop:hitchin-conductor}
Any formula for the Koszul conductor of the Hitchin quantum vertex
group must supply:
\begin{enumerate}[label=(\roman*)]
\item the weighted trace realizing
\[
\kappa_{\mathrm{ch}}(A_{C,\mathbf{G}}) =
\frac{(k + h^\vee) \cdot d_\fkg}{2 h^\vee} \cdot (g_C - 1)
\]
on the affine local model;
\item the Langlands-dual weighted trace for ${}^L\mathbf{G}$;
\item the cyclic modular Koszul pair identifying
$A_{C,\mathbf{G}}^!$ with the appropriate dual object;
\item the conductor $K_{C,\mathbf{G}}$ extracted from that pair.
\end{enumerate}
The closed expression
$(g_C - 1)d_\fkg(h^\vee+h^{\vee,L})$ is a normalization target, not a
proved formula in the absence of these data.
\end{proposition}

\begin{proof}[Criterion]
The displayed affine formula is the local Kac--Moody target. A Hitchin
conductor theorem must first prove that the Hitchin $\Phi_d$ output has
that weighted trace, then prove the dual trace and the Koszul pairing.
Only then can the conductor be evaluated.
\end{proof}

\subsection{The Feigin--Frenkel--Langlands triangle}

The three dualities (chiral Koszul, Feigin--Frenkel, and
Langlands) form a triangle:

\begin{equation}\label{eq:fff-triangle}
\begin{array}{ccc}
\widehat{\fkg}_k & \xrightarrow{\;\text{chiral Koszul}\;} &
\widehat{\fkg}_{-k-2h^\vee} \\[6pt]
\Big\downarrow{\scriptstyle\text{Langlands}} & &
\Big\downarrow{\scriptstyle\text{Langlands}} \\[6pt]
\widehat{\fkg^L}_{k^L} &
\xrightarrow{\;\text{chiral Koszul}\;} &
\widehat{\fkg^L}_{-k^L-2h^{\vee,L}}
\end{array}
\end{equation}
where $k^L$ satisfies $(k+h^\vee)(k^L + h^{\vee,L}) = 1$. The
vertical arrows are \emph{not} chiral Koszul duality; they are
Feigin--Frenkel/Langlands duality, a different operation. The
conjecture $G(C, \mathbf{G})^! = G(C, {}^L\mathbf{G})$ asserts that
\emph{at the level of quantum vertex chiral groups}, the horizontal
and vertical arrows compose to give the diagonal: the passage from
$G(C, \mathbf{G})$ to $G(C, {}^L\mathbf{G})$ is Koszul duality.

\begin{remark}[Distinction from affine Kac--Moody Koszul duality]
For affine Kac--Moody algebras, chiral Koszul duality is
$k \mapsto -k - 2h^\vee$ \emph{within the same Lie algebra} $\fkg$.
The Langlands dual algebra $\fkg^L$ appears only through the
Feigin--Frenkel center at critical level, not through Koszul duality.
The new content of Conjecture~\ref{conj:langlands-koszul} is that
for \emph{quantum vertex chiral groups} (which incorporate the full
BPS/automorphic data), Koszul duality \emph{does} exchange
$\mathbf{G}$ and ${}^L\mathbf{G}$. The extra ingredient is the
center-layer braiding and the automorphic correction (shadow
obstruction tower), which are invisible at the level of affine
Kac--Moody algebras alone.
\end{remark}


% ======================================================================
\section{The Borcherds involution and BKM Koszul duality}
\label{sec:borcherds}
% ======================================================================

\subsection{Why Cartan negation fails}

For a finite-dimensional simple Lie algebra $\fkg$ with Cartan matrix
$A = (a_{ij})$, the Koszul dual algebra is $\fkg$ itself (self-duality
of the Koszul pair $(\mathrm{Lie}, \mathrm{Com})$). The Cartan matrix
is unchanged.

For a Kac--Moody algebra with symmetrizable generalized Cartan matrix,
the situation is similar: the Koszul dual is the same Kac--Moody
algebra at a shifted level (the Feigin--Frenkel involution).

For a generalized BKM superalgebra $\fkg_\Phi$ (where $\Phi$ is an
automorphic form controlling the root multiplicities), the naive
prescription ``negate the Cartan matrix'' fails for three reasons:

\begin{enumerate}[label=(\arabic*)]
\item The generalized Cartan matrix of a BKM algebra may have
$a_{ii} \leq 0$ for some $i$ (imaginary simple roots). Negating
gives $-a_{ii} \geq 0$, which changes the type of the root
(imaginary $\leftrightarrow$ real) and hence the structure of the
algebra.

\item The imaginary root multiplicities $\mult(\alpha) = f(nm, l)$
are determined by the automorphic form $\Phi$, not by the Cartan
matrix. Negating the Cartan matrix does not specify what the dual
multiplicities should be.

\item The positivity structure (the choice of positive cone
$\cR_{\geq 0} \mathcal{P}$) is part of the BKM data. Negation
reverses the positive cone, which is not a well-defined operation
on the BKM algebra unless accompanied by a specific involution that
preserves the automorphic structure.
\end{enumerate}

\subsection{The correct prescription: Borcherds involution}

\begin{definition}[Borcherds involution for BKM superalgebras]
\label{def:borcherds-involution-general}
Let $\fkg_\Phi$ be a generalized BKM superalgebra with root lattice
$\Lambda$, Weyl group $W$, automorphic form $\Phi$, and positive
cone $C_+ \subset \Lambda \otimes \mathbb{R}$. The \emph{Borcherds
involution} $\omega_B \colon \fkg_\Phi \to \fkg_\Phi^!$ is the
unique algebra involution satisfying:
\begin{enumerate}[label=(\roman*)]
\item $\omega_B$ restricts to the Chevalley involution on real root
spaces: $\omega_B(e_\alpha) = -f_\alpha$,
$\omega_B(f_\alpha) = -e_\alpha$, $\omega_B(h_\alpha) = -h_\alpha$
for $\alpha \in \Delta^{\mathrm{re}}$.
\item $\omega_B$ permutes imaginary root spaces according to the
automorphic involution: $\omega_B(\fkg_\alpha) = \fkg_{\sigma(\alpha)}$
for $\alpha \in \Delta^{\mathrm{im}}$, where
$\sigma \colon C_+ \to C_+$ is the involution induced by the
functional equation of $\Phi$.
\item $\omega_B$ is compatible with the Weyl group:
$\omega_B \circ w = w \circ \omega_B$ for all $w \in W$.
\end{enumerate}
\end{definition}

\begin{proposition}
The Borcherds involution $\omega_B$ is well-defined and unique,
provided the automorphic form $\Phi$ has a functional equation
$\Phi(\sigma \cdot z) = (\det \sigma)^{-\kappa_{\mathrm{ch}}} \Phi(z)$ for some
involution $\sigma$ of the domain.
\end{proposition}

\begin{proof}
The functional equation of $\Phi$ ensures that the root multiplicities
satisfy $\mult(\alpha) = \mult(\sigma(\alpha))$ (up to the sign from
$(\det \sigma)^{-\kappa_{\mathrm{ch}}}$, which accounts for the shift $\rho \to \rho^!$
in the Weyl vector). The Chevalley involution on real roots is standard
(Kac). The extension to imaginary roots by $\sigma$ is uniquely
determined because $\sigma$ is an involution of the positive cone.
\end{proof}

\subsection{The dual automorphic form}

The denominator identity of $\fkg_\Phi^!$ involves the
\emph{dual automorphic form}
\begin{equation}\label{eq:dual-automorphic}
\Phi^!(z) = \Phi(\sigma \cdot z) \cdot (\det \sigma)^{-\kappa_{\mathrm{ch}}(\Phi)}.
\end{equation}

For $K3 \times E$: $\Phi = \Delta_5$, $\sigma = W_N$ (Fricke
involution), $\kappa_{\mathrm{ch}} = 5$, and
$\Phi^!(z) = \Delta_5(W_N \cdot z) \cdot N^{-10}$.

For toric CY3: $\Phi = Z^{\mathrm{DT}}(X_\Sigma)$, $\sigma$ is the
toric involution, and $\Phi^!$ is the DT partition function of the
``Langlands-dual toric variety'' (obtained by transposing the toric
diagram when $\fkg_{Q_X}$ is not self-dual).


% ======================================================================
\section{Four faces of CY Koszul duality}
\label{sec:synthesis}
% ======================================================================

The four constructions have different proof surfaces.

\begin{center}
\renewcommand{\arraystretch}{1.4}
\begin{tabular}{lcccc}
& \textbf{Toric CY3} & \textbf{$K3 \times E$}
& \textbf{Vol.~I chiral} & \textbf{Hitchin} \\
\hline
$G(X)$ & $Y(\widehat{\fkg})$ & $\fkg_{\Delta_5}$
& $A_X$ & $G(C,\mathbf{G})$ \\
$G(X)^!$ & $Y(\widehat{\fkg^L})$ & $\fkg_{\Delta_5}^!$
& $A^!_X$ & $G(C,{}^L\mathbf{G})$ \\
Mechanism & Y/W duality & Borcherds inv.
& Verdier on $\Ran$ & Langlands \\
Conductor $K$ & proof obligation & proof obligation
& $\varrho \cdot K_X$ on constructed pairs & proof obligation \\
Self-dual & $\fkg = \fkg^L$ & $X = X^\vee$
& $c = c_*$ & $\mathbf{G}$ s.l. \\
\end{tabular}
\end{center}

\noindent Abbreviations: s.l.\ = simply-laced, s.m.\ = self-mirror.

\medskip

The unifying principle is:

\begin{center}
\fbox{\parbox{0.85\textwidth}{
\textbf{Koszul duality of completed quantum vertex chiral groups is a
candidate $E_2$-center incarnation of Langlands duality.}
The passage from $G(X)$ to $G(X)^!$ replaces the structure group
$\mathbf{G}$ by ${}^L\mathbf{G}$, the level $k$ by its
Feigin--Frenkel dual, and the imaginary roots by their Borcherds
involution images, after the positive half, pairing, completion, and
recognition data are fixed.
}}
\end{center}

\subsection{Compatibility with the main theorems}

\begin{itemize}
\item \textbf{Theorem CY-A} (CY-to-chiral assignment): The
dimension-stratified assignment $\Phi_d$ produces an $E_1$ chiral output
at $d\geq 3$ after specialization. The claim
$A_{C,\mathbf{G}}^! = A_{C,{}^L\mathbf{G}}$ requires a constructed
Koszul pair.

\item \textbf{Theorem CY-B} (bar-cobar and Verdier duality): The
native bar object is $B_{E_1}(A_{C,\mathbf{G}})$. Bar-cobar inversion
recovers $A_{C,\mathbf{G}}$; Verdier duality of bar cohomology produces
$A_{C,\mathbf{G}}^!$.

\item \textbf{Theorem CY-C} (quantum group realization): The braided
category lives on the center or completed double. A Langlands-dual
braided equivalence is a theorem only on loci where both doubles and
the braiding reversal $q \mapsto q^{-1}$ have been constructed.

\item \textbf{Theorem CY-D} (modular characteristic):
$\kappa_{\mathrm{ch}}(A_{C,\mathbf{G}}) + \kappa_{\mathrm{ch}}(A_{C,{}^L\mathbf{G}}) =
\varrho \cdot K_{C,\mathbf{G}}$, where
$K_{C,\mathbf{G}}$ is the conductor of the constructed Koszul pair.
\end{itemize}


% ======================================================================
\section{Open questions}
\label{sec:open}
% ======================================================================

\begin{enumerate}[label=\textbf{Q\arabic*.}]
\item \textbf{Non-toric, non-$K3 \times E$ CY3.}
For the quintic threefold $X_5 \subset \bP^4$, or for a generic CY3
with $h^{2,1} > 0$: what is $G(X_5)^!$? Is there a Langlands-type
duality for CY3 that are not associated to any gauge group?

\item \textbf{Wall-crossing and Koszul duality.}
The Kontsevich--Soibelman wall-crossing formula changes the BPS
spectrum $\{\DT_\alpha(X)\}$ as one crosses walls of marginal
stability. How does Koszul duality interact with wall-crossing?
Conjecture: wall-crossing is a gauge equivalence of MC elements
$\Theta_{A_X} \sim \Theta'_{A_X}$, and Koszul duality commutes with
gauge equivalence.

\item \textbf{Higher-genus extension.}
The Koszul dual denominator identity
(Proposition~\ref{prop:k3e-koszul-denominator}) lives on $\bH_2$
(genus~2). Is there a genus-$g$ Koszul duality that relates
denominator identities on the genus-$g$ Siegel space $\bH_g$?

\item \textbf{Quantum geometric Langlands.}
The quantum geometric Langlands conjecture
(Feigin--Frenkel--Stoyanovsky, Gaitsgory) predicts an equivalence
of categories of twisted $D$-modules on $\mathrm{Bun}_\mathbf{G}$
and $\mathrm{Bun}_{{}^L\mathbf{G}}$. Can this equivalence be
realized as Koszul duality of the associated quantum vertex chiral
groups, with the $E_2$ structure carried by the KZ center layer of the
KM algebras $V_k(\frakg)$?

\item \textbf{The eight diagonal divisor forms.}
Conjecture~\ref{conj:k3e-mirror-koszul}(iii) predicts that the
eight quantum vertex chiral groups $G(X_N)$ are paired by Koszul
duality. What is the explicit Fricke pairing? Which pairs are
self-dual, and which are genuinely dual to a different $N$?

\item \textbf{S-duality and Koszul duality.}
In the physics literature, S-duality of 4d $\cN = 4$ gauge theory
exchanges $\mathbf{G}$ and ${}^L\mathbf{G}$ with
$\tau \mapsto -1/\tau$. The reduction to 2d on $C$ gives geometric
Langlands. Is the Koszul duality
$G(C, \mathbf{G})^! = G(C, {}^L\mathbf{G})$ the 2d shadow of 4d
S-duality via the Kapustin--Witten twist?
\end{enumerate}


% ======================================================================
\section{Appendix: $\bC^3$ root data}
\label{sec:appendix-c3}
% ======================================================================

The $\bC^3$ computation separates the positive half from the
double/center layer.

\subsection{Root datum}
The native Hall algebra is the positive half
$Y^+(\widehat{\mathfrak{gl}}_1)$. Its completed Drinfeld double is the
quantum vertex group layer compared with $\mathcal W_{1+\infty}$. The
positive-half root datum has:
\begin{itemize}
\item Lattice: $\Lambda = \bZ$ (one-dimensional, from the Jordan quiver).
\item Real roots: none (the Jordan quiver has a single vertex with a
single loop; there are no finite-type roots).
\item Imaginary roots: $\Delta^{\mathrm{im}}_0 = \{n \in \bZ_{>0}\}$.
The MacMahon product
$M(q)=\prod_{n\geq 1}(1-q^n)^{-n}$ gives exponent $n$ in degree $n$;
partition numbers appear as coefficients of $M(q)$, not as the
primitive root multiplicities.
\item Weyl group: trivial.
\item Weyl vector: $\rho = 0$.
\end{itemize}

\subsection{Koszul dual root datum}
$\Lambda^! = \Lambda^\vee = \bZ$. The Borcherds involution acts as the
identity on $\Delta^{\mathrm{im}}_0$ because
$\mathfrak{gl}_1^L=\mathfrak{gl}_1$. Hence the positive-half root datum
is self-dual; the full $\mathcal W_{1+\infty}$ comparison is made only
after the double/center passage.

\subsection{Shadow obstruction tower}
The compact Hodge-supertrace formula does not apply to non-compact
$\bC^3$. In the standard self-dual $\mathcal W_{1+\infty}$ normalization
the degree-$2$ curvature vanishes, but the MacMahon product has
infinitely many nonzero exponents. The leading shadow can vanish while
higher shadows remain nonzero.

\subsection{DT partition function}
The DT partition function of $\bC^3$ is the MacMahon function
$Z^{\DT}(\bC^3)=M(q)$. The trivial Borcherds involution leaves this
product fixed. This verifies product-level self-duality; identifying it
with chiral Koszul self-duality requires the completed double and
center comparison.


\end{document}
